\documentclass[10pt,twoside]{article}
\usepackage[utf8]{inputenc}
\usepackage[spanish]{babel}
\usepackage{amsmath, amsfonts, amssymb}
\usepackage{lmodern} 
\usepackage{graphicx}
\usepackage{tikz}
\usetikzlibrary{babel, 3d, perspective}
\usepackage{tikz-3dplot} 
\usepackage{geometry}
\usepackage{xcolor}
\usepackage{fancyhdr}
\usepackage{multicol}
\usepackage{hyperref}

% Configuración de márgenes
\geometry{
  a4paper,
  top=1.5cm,
  bottom=1.5cm,
  inner=1.5cm,
  outer=1.5cm,
  headheight=15pt
}

% Paleta de colores UNAB
\definecolor{unabBlue}{HTML}{003865}
\definecolor{unabOrange}{HTML}{E85C0B}

% Encabezados
\pagestyle{fancy}
\fancyhf{}
\fancyhead[LE,RO]{\thepage}
\fancyhead[RE]{\textcolor{unabBlue}{\textit{Curso de Electromagnetismo}}}
\fancyhead[LO]{\textcolor{unabBlue}{\textit{Quiz de Aplicación}}}
\fancyfoot[C]{\textcolor{unabBlue}{\footnotesize Universidad Autónoma de Bucaramanga -- Departamento de Ciencias Básicas}}
\renewcommand{\headrulewidth}{0.4pt}
\renewcommand{\footrulewidth}{0.4pt}

% Comando de sección
\newcommand{\mysection}[1]{
  \vspace{0.3cm}\noindent{\Large\bfseries\color{unabBlue} #1}
  \vspace{0.1cm}\hrule\vspace{0.2cm}
}

% Comando para problemas
\newcounter{problema}
\newcommand{\problema}[1]{
    \stepcounter{problema}
    \vspace{0.2cm}\noindent\textbf{\textcolor{unabBlue}{Problema \theproblema}}\ (#1)\par\vspace{0.1cm}\noindent
}

\newcommand{\class}{\textcolor{unabBlue}{Tema: Electrostática y Campo Eléctrico}}
\newcommand{\term}{Semestre en Curso}
\newcommand{\examnum}{Quiz de Aplicación (RAEs)}
\newcommand{\examdate}{Fecha:}
\newcommand{\timelimit}{2 Horas}

\begin{document}

\noindent
\begin{minipage}[c]{0.12\textwidth}
    \includegraphics[width=\linewidth]{../../Plantilla_LaTeX_Examen_Ciencia_de_Datos_UNAB/logo_ciencia_datos_negro.png}
\end{minipage}%
\hspace{0.03\textwidth}%
\begin{minipage}[c]{0.85\textwidth}
\begin{tabular}{@{}p{0.55\textwidth} p{0.42\textwidth}@{}}
\textbf{\class} & \textbf{Nombre:} \hrulefill \\
\textbf{\term} & \textbf{Curso:} Electromagnetismo \\
\textbf{\examnum} & \textbf{Fecha:} \hrulefill \\
\textbf{Tiempo: \timelimit} & \textbf{Profesor:} Dudbil Olvasada Pabon R.
\end{tabular}
\end{minipage}
\vspace{0.1cm}
\rule[2ex]{\textwidth}{1pt}

\vspace{0.1cm}
\noindent
\textbf{Instrucciones:} Este examen evalúa los conceptos teóricos y problemas analíticos de fuerza electrostática, campo y potencial. Resuelva cada punto de forma clara. \textbf{No se permite el uso de celulares ni formular preguntas teóricas durante la prueba. El incumplimiento de estas normas dará lugar a sanciones disciplinarias.}

\vspace{0.2cm}
\begin{multicols}{2}

\mysection{I. Análisis Conceptual y Teórico}

\problema{Valor: 1.0 punto | Evalúa RAE 1.1, 1.2, 2.1, 2.2}
Defina de forma clara y rigurosa los siguientes conceptos fundamentales de la electrostática:
\begin{itemize}
    \item \textbf{Carga Eléctrica:} ¿Qué es desde un punto de vista físico fundamental?
    \item \textbf{Tipos de Electrización:} Describa brevemente los tres mecanismos físicos para electrizar un material (fricción, inducción y contacto).
    \item \textbf{Campo Eléctrico y Potencial:} Defina el vector Campo Eléctrico y diferéncielo conceptualmente del Potencial Eléctrico.
    \item \textbf{Líneas y Superficies:} ¿Qué representan las líneas de campo eléctrico y qué relación geométrica estricta deben cumplir respecto a las superficies equipotenciales?
\end{itemize}

\vspace{0.2cm}

\mysection{II. Modelado Matemático y Resolución de Problemas}

\problema{Valor: 1.5 puntos | Evalúa RAE 1.3 y 2.1}
Tres cargas puntuales se encuentran fijas en los vértices de un triángulo rectángulo. La carga $q_1 = -4\ \mu\text{C}$ está en el origen $(0,0)$. La carga $q_2 = +5\ \mu\text{C}$ está ubicada en $(0, 3)\text{ m}$, y la carga $q_3 = +3\ \mu\text{C}$ está en $(4, 0)\text{ m}$. 
\textbf{a)} Calcule de forma vectorial la \textbf{Fuerza Electrostática neta} que actúa sobre $q_1$. 
\textbf{b)} Calcule el \textbf{Campo Eléctrico neto} en el origen producido por $q_2$ y $q_3$. 
(Para ambos, indique magnitud total y ángulo respecto al eje $X$).

\vspace{0.2cm}

\problema{Valor: 1.0 punto | Evalúa RAE 2.1}
En los vértices de un cuadrado de lado $L = 2\text{ m}$ en el plano $XY$ se ubican tres cargas puntuales: $q_A = +2\ \mu\text{C}$ en $(0,2)\text{ m}$, $q_B = -4\ \mu\text{C}$ en $(2,2)\text{ m}$ y $q_C = +5\ \mu\text{C}$ en $(2,0)\text{ m}$.
Calcule el \textbf{Potencial Eléctrico neto} ($V_{\text{neto}}$) evaluado en el origen $O(0,0)$ y determine el \textbf{Trabajo} eléctrico externo necesario para traer una partícula de carga $q = +1\ \mu\text{C}$ desde el infinito hasta dicho origen.

\vspace{0.2cm}

\problema{Valor: 1.5 puntos | Evalúa RAE 2.2}
Una varilla delgada aislante de longitud $L = 2\text{ m}$ se encuentra alineada a lo largo del eje $X$, con su extremo izquierdo en el origen $(0,0)$. La varilla posee una densidad de carga lineal uniforme $\lambda = +12\ \text{nC/m}$.
Planteando desde los diferenciales geométricos, determine la magnitud y dirección del \textbf{campo eléctrico} en un punto $P$ ubicado sobre el eje $X$ a una distancia $a = 1\text{ m}$ a la izquierda del origen (es decir, en la coordenada $x = -1\text{ m}$). Justifique matemáticamente la integración.

\mysection{III. Apéndice y Formulario}
\vspace{0.3cm}
\textbf{Constantes Físicas:}
\begin{itemize}
    \item $k_e = \frac{1}{4\pi\epsilon_0} \approx 8.99 \times 10^9\ \text{N}\cdot\text{m}^2/\text{C}^2$
    \item $\epsilon_0 = 8.85 \times 10^{-12}\ \text{C}^2/(\text{N}\cdot\text{m}^2)$
    \item $1\ \mu\text{C} = 10^{-6}\ \text{C}$
    \item $1\ \text{nC} = 10^{-9}\ \text{C}$
\end{itemize}

\textbf{Magnitud y Vectores:}
\begin{itemize}
    \item Vector posición: $\vec{r} = x\hat{i} + y\hat{j} + z\hat{k}$
    \item Magnitud: $|\vec{r}| = \sqrt{x^2 + y^2 + z^2}$
    \item Vector Desplazamiento: $\vec{R} = \vec{r}_{\text{pto}} - \vec{r}_{\text{fuente}}$
    \item Vector Unitario: $\hat{u} = \frac{\vec{R}}{|\vec{R}|}$
\end{itemize}

\textbf{Electrostática Discreta:}
\begin{itemize}
    \item Fuerza de Coulomb: $\vec{F}_{12} = k_e \frac{q_1 q_2}{|\vec{R}|^2} \hat{u}_{R}$
    \item Campo Eléctrico: $\vec{E} = k_e \frac{q}{|\vec{R}|^2} \hat{u}_{R}$
    \item Fuerza y Campo: $\vec{F} = q_0 \vec{E}$
    \item Potencial: $V = k_e \frac{q}{|\vec{R}|}$
\end{itemize}

\textbf{Distribuciones Continuas:}
\begin{itemize}
    \item Lineal: $dq = \lambda\ dl$
    \item Superficial: $dq = \sigma\ dA$
    \item Volumétrica: $dq = \rho\ dV$
    \item Campo: $\vec{E} = \int k_e \frac{dq}{|\vec{R}|^2} \hat{u}_{R}$
    \item Potencial: $V = \int k_e \frac{dq}{|\vec{R}|}$
\end{itemize}

\end{multicols}
\end{document}
